\chapter{Технологический раздел}


\section{Описание входных и выходных данных}


\section{Метрики оценки качества реконструкции глубины}

\subsection{Абсолютная средняя ошибка (MAE)}
\begin{equation}
\text{MAE} = \frac{1}{n} \sum_{i=1}^{n} |d_{i}^{\text{pred}} - d_{i}^{\text{gt}}|
\end{equation}

MAE представляет собой среднее абсолютное отклонение предсказанной глубины от эталонной. Эта метрика менее чувствительна к выбросам по сравнению с RMSE и позволяет выявлять систематические ошибки в предсказании глубины~\cite{zhang2025review}.

\subsection{Среднеквадратичная ошибка (RMSE)}
\begin{equation}
\text{RMSE} = \sqrt{\frac{1}{n} \sum_{i=1}^{n} (d_{i}^{\text{pred}} - d_{i}^{\text{gt}})^2}
\end{equation}

RMSE придаёт больший вес крупным ошибкам и является стандартной метрикой для оценки точности методов оценки глубины~\cite{vis3}. Благодаря своей чувствительности к выбросам, данная метрика эффективна при оценке методов, где критично минимизировать значительные отклонения.

\subsection{Относительная абсолютная ошибка (REL)}
\begin{equation}
\text{REL} = \frac{1}{n} \sum_{i=1}^{n} \frac{|d_{i}^{\text{pred}} - d_{i}^{\text{gt}}|}{d_{i}^{\text{gt}}}
\end{equation}

REL учитывает относительную величину ошибки, что особенно важно при работе с объектами, находящимися на разном расстоянии от камеры. Эта метрика предпочтительна при оценке качества реконструкции исторических объектов, где точная ш��ала глубины может быть неизвестна~\cite{zhang2025review}.

\subsection{Точность $\delta$}
\begin{equation}
\delta_{k} = \frac{1}{n} \sum_{i=1}^{n} \mathbb{I}\left[\max\left(\frac{d_{i}^{\text{pred}}}{d_{i}^{\text{gt}}}, \frac{d_{i}^{\text{gt}}}{d_{i}^{\text{pred}}}\right) < 1.25^k\right]
\end{equation}

где $k \in \{1, 2, 3\}$, а $\mathbb{I}[\cdot]$ — индикаторная функция, равная 1, если условие выполняется, и 0 в противном случае.

Эта метрика определяет процент пикселей, для которых отношение предсказанной глубины к истинной (или наоборот) не превышает заданного порога. Стандартно используются три порога: $1.25$, $1.25^2$ и $1.25^3$~\cite{vis2}.

\subsection{Перцептивное качество}
Субъективная оценка реалистичности полученной 3D-модели особенно важна при работе с историческими фотографиями, где эталонные данные недоступны.


В контексте реконструкции глубины из архивных фотографий особое внимание следует уделять метрикам без эталонных данных, так как получить достоверную эталонную информацию о глубине для исторических сцен практически невозможно~\cite{zaharov2011}.


\section{Тестирование}

\subsection{Подход к тестированию}

В контексте оценки методов реконструкции глубины из исторических фотографий возникает фундаментальная проблема: отсутствие эталонных данных о глубине для архивных снимков.
Для решения этой проблемы предлагается подход, основанный на искусственном старении современных изображений с известной картой глубины.
Такой метод позволяет:

\begin{itemize}
    \item Использовать объективные количественные метрики (MAE, RMSE, REL, $\delta_k$), широко признанные в научном сообществе~\cite{vis3}
    \item Контролировать степень и характер эффектов старения
    \item Проводить систематическое тестирование с изолированием отдельных факторов деградации изображений
    \item Получать статистически значимые результаты для различных типов сцен
\end{itemize}

\subsection{Классы эквивалентности и граничные условия}

Для систематического тестирования метода реконструкции глубины были выделены следующие классы эквивалентности:

\begin{table}[H]
\centering
\caption{Классы эквивалентности для тестирования}
\label{tab:equivalence-classes}
\begin{tabular}{|p{3cm}|p{9cm}|}
\hline
\textbf{Критерий классификации} & \textbf{Классы эквивалентности} \\
\hline
Тип сцены &
\begin{itemize}
    \item Архитектурные объекты (отдельные здания)
    \item Городские панорамы (многоплановые сцены)
    \item Групповые портреты (люди в интерьере)
    \item Интерьерные сцены
    \item Пейзажные снимки
\end{itemize} \\
\hline
Исторический период имитации &
\begin{itemize}
    \item Ранние фотографии (1850-1900)
    \item Довоенные фотографии (1900-1939)
    \item Военный и послевоенный период (1939-1960)
    \item Поздние аналоговые фотографии (1960-1990)
\end{itemize} \\
\hline
Типы эффектов старения &
\begin{itemize}
    \item Механические повреждения (царапины, изломы)
    \item Химические дефекты (пятна, выцветание)
    \item Оптические дефекты (размытие, засветка)
    \item Комбинированные эффекты
    \item Изображения без эффектов (контрольная группа)
\end{itemize} \\
\hline
Степень деградации &
\begin{itemize}
    \item Минимальная (едва заметные эффекты)
    \item Умеренная (типичная для хорошо сохранившихся фотографий)
    \item Значительная (выраженные эффекты старения)
    \item Экстремальная (сильные повреждения, тест на отказоустойчивость)
\end{itemize} \\
\hline
Источник эталонных данных &
\begin{itemize}
    \item Датчики глубины (NYU Depth v2, ...)
    \item Стереоизображения (ETH3D, ...)
    \item Лазерное сканирование (Make3D, ...)
    \item Синтетические данные (Blender, Unity)
\end{itemize} \\
\hline
\end{tabular}
\end{table}

Граничные условия для тестирования включают:
\begin{itemize}
    \item Экстремальная степень деградации изображения при сохранении распознаваемости сцены
    \item Сочетание различных типов эффектов старения в одном изображении
    \item Пограничные случаи оптических эффектов (сильное размытие, экстремальный контраст)
    \item Исторические оптические искажения, характерные для старинных камер
\end{itemize}

\subsection{Тестовый набор данных}

Тестовый набор данных включает три основные категории:

\subsubsection{Современные изображения с искусственным старением}

Основу тестового набора составляют современные изображения с известной глубиной, подвергнутые различным эффектам старения:

\begin{itemize}
    \item \textbf{NYU Depth v2} --- 50 интерьерных сцен с эталонной глубиной от Kinect, подвергнутых эффектам старения разных исторических периодов
    \item \textbf{ETH3D} --- 35 архитектурных сцен с прецизионной глубиной, с применением эффектов старения разной интенсивности
    \item \textbf{Make3D} --- 25 пейзажных и городских сцен с картами глубины, обработанные для имитации фотографий 1850-1990 годов
\end{itemize}

Для каждого изображения создаются варианты с различной степенью деградации и комбинациями эффектов, что позволяет оценить влияние конкретных факторов на точность реконструкции глубины.

\subsubsection{Синтетические тестовые данные}

\begin{itemize}
    \item \textbf{Модели исторических зданий в Blender} --- 15 трехмерных моделей архитектурных памятников с рендерингом в стиле фотографий разных исторических периодов
    \item \textbf{Городские сцены из Unity} --- 10 детализированных городских сцен с точной глубиной и реалистичным рендерингом, имитирующим исторические фотопроцессы~\cite{istmin2024}
\end{itemize}

Синтетические данные обеспечивают абсолютную точность эталонных карт глубины и позволяют создавать сцены, идентичные историческим фотографиям.

\subsubsection{Контрольный набор архивных фотографий}

\begin{itemize}
    \item \textbf{Коллекция Shorpy} --- 25 высококачественных оцифрованных архивных фотографий
    \item \textbf{Архив Library of Congress} --- 15 архитектурных фотографий из коллекции HABS
    \item \textbf{Региональные архивы} --- 15 исторических снимков местных архитектурных памятников
\end{itemize}

Хотя для архивных фотографий отсутствует эталонная глубина, они используются для качественной оценки и визуального сравнения с результатами тестирования на искусственно состаренных изображениях.

\subsection{Методика тестирования}

\subsubsection{Процесс создания искусственно состаренных изображений}

Для объективной оценки влияния эффектов старения на качество реконструкции глубины используется следующий процесс:

\begin{enumerate}
    \item Выбор современного изображения с известной картой глубины
    \item Применение набора эффектов старения с параметрами, соответствующими конкретному историческому периоду:
    \begin{itemize}
        \item Имитация зернистости фотоматериалов соответствующей эпохи
        \item Добавление оптических искажений, характерных для исторических объективов
        \item Наложение эффектов химической деградации (выцветание, изменение контраста)
        \item Добавление механических повреждений (царапины, сгибы, пятна)
    \end{itemize}
    \item Создание нескольких вариантов с градацией интенсивности эффектов
    \item Реконструкция глубины из обработанного изображения
    \item Количественное сравнение с эталонной глубиной
\end{enumerate}

Для имитации эффектов старения используются как существующие библиотеки обработки изображений, так и специально разработанные алгоритмы, воспроизводящие физические и химические процессы деградации фотографий~\cite{zaharov2014}.

\subsubsection{Типы проводимых тестов}

\begin{itemize}
    \item \textbf{Тесты базовой точности} --- сравнение точности реконструкции на необработанных современных изображениях и на изображениях с минимальными эффектами старения

    \item \textbf{Тесты устойчивости к деградации} --- оценка изменения точности с увеличением интенсивности эффектов старения

    \item \textbf{Визуальная валидация} --- сопоставление результатов, полученных на искусственно состаренных изображениях, с результатами для реальных исторических фотографий
\end{itemize}

\subsection{Результаты тестирования}

Результаты тестирования представлены в таблице~\ref{tab:test-results}. Для каждого тестового случая приводятся количественные метрики на основе эталонных данных.

\section{}



